\FloatBarrier
\section{DC region-1 back membranes}
\label{sec:dc-back}

Upcoming in GEMC2 6.2 and CLAS12 GEMC Systems 0.4, the six dcs1r1back through dcs6r1back Kapton membranes move
from magnets to the DC system. A new DC cad directory owns one unchanged dcs5r1back STL. The other five named
placements share its solid. GEMC3 publishes all six passive placements in dc\_cad; GEMC2 uses one CAD source
in dc\_cad and five native CopyOf rows in dc\_copies. The supplied DC cards and the GEMC3 experiment card load
the required companions. Native DC geometry and materials remain unchanged.

The six engineering files are not byte-identical: sector coordinates and small tilts are baked into the
vertices. Proper rigid fits of their twelve corresponding vertices give a maximum error of 0.00007613 mm
against the shared reference. This is comparable to single-precision STL export rounding. Sector-specific
translations are retained rather than assuming exact rotations of a single unshifted sector. Original meshes
are archived separately in the owning GEMC2 source repository.

Each source has twenty triangles, twelve vertices, thirty edges, one closed manifold component, genus zero,
and no boundary edges. These hexagonal prisms already have the minimum closed triangular surface for their
twelve retained corners. No geometric decimation or hole filling is applied. Unique stored and built triangles
fall from 120 to 20 (83.33\% fewer), and active STL storage falls from 6,504 to 1,084 bytes. The six physical
membranes still have 120 placed triangles. Magnets now has 45 STL files with 126,658 triangles; DC owns the
additional twenty unique triangles.

Sharing the rounded approximately 2.5-micrometre reference film changes total signed Kapton volume from
23637.244828 to 23756.565606 cubic millimetres (+0.504800\%). Float coordinate rounding is significant
relative to this thickness. This numerical material difference is reported separately from the storage
reduction; it does not establish detector-response equivalence.

At the user's request, angles are rounded to integers after comparing them with the fitted angles. The largest
resulting vertex displacement is 0.759079 mm in sector 1; sectors 2 and 3 move by at most 0.634661 mm. The
comparison loads native DC, both DC and magnets CAD, the solenoid, native beamline, and beamline CAD at run
11: 139 volumes including the world. Both variants use the same seed and 100,000 Geant4 surface samples per
tested volume. No overlap involves any of the six membranes in either variant, and both variants report 48
overlap messages covering the same 23 distinct ordered volume pairs. Existing overlaps elsewhere are retained.
Surface sampling is not a mathematical proof of clearance.

All three inherited CAD runs (11, 3029, 4763) pass GEMC2 TEXT/CAD, GEMC2 SQLite/CAD, and GEMC3 initialization.
GEMC3 also constructs its custom DDVCS run, for ten successful checks in total. The semantic comparison
matches both native DC variations and all four DC CAD banks, along with all three reduced magnets CAD banks.
Independent database checks verify the GEMC2 five-copy and one-mesh rows, unchanged unrelated native/material
banks, byte-identical source/runtime meshes in both repositories, and the six unchanged archived originals.

The timing fixture isolates the six membranes at run 11. Before optimization it builds six separate
tessellated solids; afterward it builds one shared solid with five copies, using the final integer angles. One
fixed GEMC3 dev executable and Geant4 11.4.3 run FTFP\_BERT with default production cuts, one thread, no
field, no optical physics, no GUI, no event output, and no target or pileup. Each 5 GeV/c electron, proton,
positive pion, or neutron starts at the origin, uniformly in solid angle from 5 to 45 degrees with full
azimuth. Five alternating-order repetitions use 2,000 matched LUND primaries per particle and seeds
12345--12349. Startup and beamOn tracking are reported separately; forty tracking processes each report 2,000
events. All fifty timed processes and two separate construction checks complete without Geant4 exceptions.

Median zero-event process wall time falls from 0.72316 to 0.40457 s (1.787 times faster). Median process CPU
time changes from 0.660 to 0.360 s and peak RSS from 118.45 to 73.34 MiB. Tracking ratios below are
before/after; values above one indicate faster tracking. They describe this thin-membrane workload and do not
establish a full-CLAS12 speedup. The physical surface change from angle rounding and the small material-volume
difference are included in these timings.

\begin{table}[htbp]
\centering
\setlength{\tabcolsep}{8pt}
\caption{DC membrane tracking; five repetitions of 2,000 events per particle.}
\label{tab:dc-back}
\begin{tabular}{lrrrr}
\toprule
Particle & Before (s) & After (s) & Speedup & Paired ratio range \\
\midrule
Electron & 2.18817 & 2.22178 & 0.985 & 0.963--1.002 \\
Proton & 3.31635 & 3.36257 & 0.986 & 0.972--1.012 \\
Positive pion & 3.04980 & 3.08063 & 0.990 & 0.970--1.007 \\
Neutron & 0.91443 & 0.91720 & 0.997 & 0.965--0.999 \\
\bottomrule
\end{tabular}
\end{table}

\begin{figure}[htbp]
\centering
\includegraphics[width=\linewidth]{figures/dc_back_comparison.pdf}
\caption{
World x-y projections of the six original membranes and their shared-mesh copies, with sector numbers
retained. The right panel shows the maximum vertex displacement caused by the accepted integer-angle
placement, including the small rigid-fit residual. Lines overlap at the displayed millimetre scale; the bar
chart exposes the submillimetre changes.}
\label{fig:dc-back}
\end{figure}

data/dc\_back and dc\_back\_archive.zip retain shape fits, rounded-angle displacements, database checks, both
dense assembly checks, repeated timing inputs/results, and raw logs. The archive includes the original and
final fixtures, archived engineering meshes, builder and validation scripts, and executable hash. The
executable itself is not redistributed. Earlier study archives are unchanged. Rebuild the shared mesh with the
builder's --round-angles option.
